The Inter-Integrated Circuit (I2C or I$^2$C) peripheral is a high speed, two-wire, synchronous, multi-drop serial interface. It transfers one bit at a time in a single direction. The I2C bus requires two wires, SDA and SCL. The main features of the I2C peripheral are listed below:

\begin{itemize}
	\item Master or slave mode, each with transmitter or receiver mode
	\item 8-bit transfer lengths
	\item Programmable transfer rate
	\item Continuous transfer operation to reduce dead time
	\item Interrupts for: start condition received, stop condition received, master/slave transfer complete, master/slave NACK received, master/slave transmit register empty, master arbitration lost, start condition sent, stop condition sent, slave receive register overflow, and slave addressed events
	\item Automatic master NACK before stop or restart condition
	\item Optional slave clock stretching
	\item 7-bit slave address with optional address bit mask
	\item Optional slave general call ability
\end{itemize}

A summary of the I2C protocol can be found at {\color{blue}\href{https://www.nxp.com/docs/en/user-guide/UM10204.pdf}{\ul{https://www.nxp.com/docs/en/user-guide/UM10204.pdf}}}


\subsection{I2C Pins}

\begin{figure}[h]
	\includesvg{figures/i2c-connection-diagram}
	\caption{I2C Connection Diagram}
	\label{fig:i2c-connection-diagram}
\end{figure}

The I2C bus has two pins: \pin{SDAx} and \pin{SCLx}. A connection diagram for the I2C bus is shown in Figure \ref{fig:i2c-connection-diagram}.

\pin{SDAx} is the serial data pin, which transmits each bit. It must be attached to the SDA pin of every other device that shares the I2C bus, and must have a pullup resistor to $V_{DD}$ attached to the line, preferably externally. When the corresponding \register{PxSEL[y]} bit is a 1, the \peripheral{I2Cx} peripheral seizes the \texttt{DIR} and \texttt{OUT} controls of \pin{SDAx}, but the \texttt{REN} controls are always configurable by the corresponding \register{PxREN[y]} bit.

\pin{SCLx} is the serial clock pin, which governs when each transmitted bit is updated and latched. It must be attached to the SCL pin of every other device that shares the I2C bus, and must have a pullup resistor to $V_{DD}$ attached to the line, preferably externally. When the corresponding \register{PxSEL[y]} bit is a 1, the \peripheral{I2Cx} peripheral seizes the \texttt{DIR} and \texttt{OUT} controls of \pin{SCLx}, but the \texttt{REN} controls are always configurable by the corresponding \register{PxREN[y]} bit.

When the \peripheral{I2Cx} peripheral has control over the \pin{SDAx} and \pin{SCLx} pins, the pins operate in open drain (aka open collector) mode. In this mode, the \peripheral{I2Cx} peripheral can only strongly enforce a logic low signal (aka driving the pin low), but cannot strongly enforce a logic high signal (aka driving the pin high). Instead, when the \peripheral{I2Cx} peripheral wants to send a logic high signal, it releases the pins from strong enforcement, allowing the external pullup resistor to passively pull the SDA or SCL wire high. Conversely, when the \peripheral{I2Cx} peripheral wants to send a logic low signal, it drives the pin low, overpowering the external pullup resistor to strongly enforce a logic low level on the wire. This is done so that it is not possible for two devices to simultaneously attempt to drive the same wire to opposite logic levels (both low and high), which would result in nondeterministic, erroneous behavior and a large power draw.



\subsection{I2C Transfer Process}

I2C is a multidrop bus, meaning that the SDA or SCL wire will be driven to a logic low level if \textit{any} devices on the bus drive it low. On the other hand, the wire will be driven to a logic high level only if \textit{all} devices on the bus release the wire (i.e. no devices are driving the wire low), allowing the external pullup resistor to weakly pull the wire high. The I2C transfer process leverages this property to prevent transmission errors and keep the number of wires to only two.

\subsubsection{Start, Restart, and Stop Conditions}

\begin{figure}[h]
	\centering
	\begin{subfigure}[b]{.25\linewidth}
		\centering
		\begin{tikztimingtable}[timing/font=\rmfamily, timing/d/text/.append style={font=\rmfamily}, xscale=1.5, yscale=2, timing/slope=0.2, timing/dslope=0.2, timing/d/background/.style={fill=white}]
			SDA & 2H 3L \\
			SCL & 3H 2L \\
		\end{tikztimingtable}
		\caption{I2C Start Condition}
		\label{fig:i2c-condition-start}
	\end{subfigure}
	\begin{subfigure}[b]{.4\linewidth}
		\centering
		\begin{tikztimingtable}[timing/font=\rmfamily, timing/d/text/.append style={font=\rmfamily}, xscale=1.5, yscale=2, timing/slope=0.2, timing/dslope=0.2, timing/d/background/.style={fill=white}]
			SDA & ;[dotted] D{}; 3D{N/ACK} 2H 2L ;[dotted] L; \\
			SCL & ;[dotted] L; L H 2L 2H L ;[dotted] L; \\
		\end{tikztimingtable}
		\caption{I2C Restart Condition}
		\label{fig:i2c-condition-restart}
	\end{subfigure}
	\begin{subfigure}[b]{.25\linewidth}
		\centering
		\begin{tikztimingtable}[timing/font=\rmfamily, timing/d/text/.append style={font=\rmfamily}, xscale=1.5, yscale=2, timing/slope=0.2, timing/dslope=0.2, timing/d/background/.style={fill=white}]
			SDA & 3L 2H \\
			SCL & 2L 3H \\
		\end{tikztimingtable}
		\caption{I2C Stop Condition}
		\label{fig:i2c-condition-stop}
	\end{subfigure}
	\caption{I2C Start, Restart, and Stop Conditions}
	\label{fig:i2c-conditions}
\end{figure}

If no device is communicating on the I2C bus, the bus is idle. When idle, both the SDA and SCL wires are at a logic high level, meaning that no device on the bus is driving them low and the external pullup resistors are pulling them high. The bus will stay this way until a device decides to seize control of the bus and become the master.

An I2C data transfer always begins when a device wishes to become the master and issues a start condition, shown in Figure~\ref{fig:i2c-condition-start}. To issue a start condition, the device drives the SDA line low while the SCL line is high. That device becomes the new master. This signals to all other devices on the bus that a new master has seized control of the bus, and that no other device may have control of the bus until the new master relinquishes control. Note that a start condition occurs any time there is a falling edge of SDA while SCL is high. If one occurs during an ongoing an incomplete transfer, that transfer will immediately cease, and possibly be void. When a new master seizes control of the bus by issuing a start condition, it becomes responsible for also issuing all of the clock pulses on SCL. The data on SDA is latched with every rising edge of SCL, and SDA can only be updated while SCL is low, otherwise an unintentional start or stop condition will be generated.

At the end of a transfer, the master must issue a stop condition, shown in Figure~\ref{fig:i2c-condition-stop}. To issue a stop condition, the master releases the SDA wire (allows it to go high) while the SCL wire is high. A stop condition tells all other devices on the bus that the master has released control of the bus and has ceased to be the master, allowing any other device to be the new master. After a stop condition, the I2C bus goes back to the idle state where both SDA and SCL are pulled high by thier external pullup resistors.

\subsubsection{Address Frame}

An I2C transfer always begins when the master issues a start (or restart) condition followed by an address frame. The address frame is composed of a 7-bit slave address, a read/write bit, and an acknowledge bit. Each bit is clocked one at a time by the master using the SCL line, latching SDA to get the transferred bit. Figure~\ref{fig:i2c-address-frame-timing-diagram} depicts an I2C address frame timing diagram. Note that the diagram is colored by which device is controlling the SDA and SCL wires, which changes at various points throughout the transfer.

The slave address is unique to the device that the master wishes to communicate with. When the master seizes control of the bus with a start condition, all other devices on the bus listen for the slave address the the master issues. The master then transmits each bit of the 7-bit slave address one at a time beginning with the MSB. If a device's address matches the slave address issued by the master, then that device becomes the slave and the others do not participate in the remainder of the transfer.

After the 7-bit slave address, the master then transmits the read/write bit, which indicates which direction data will be travelling for the remainder of the transfer. If the read/write bit is a 0, then the rest of the transfer is in write mode, and the master will transmit data to the receiving slave. If the read/write bit is a 1, the rest of the transfer is in read mode, and the slave will transmit data to the receiving master. The read/write mode does not change for the entire duration of the transfer, which ends with either a stop or a restart condition.

Finally, after the slave address and the read/write bit have been transmitted, the master momentarily releases the SDA wire to a logic high level so the slave can transmit the acknowledge bit. If the slave realizes that it has been addressed by the master and is ready to participate in the transfer, then it momentarily seizes control of the SDA line by driving it low to send an ACK (acknowledge). The master then issues a clock pulse on SCL. However, if for any reason the slave is not able to communicate with the master, then it issues a NACK (not acknowledge) by \textit{not} driving SDA low and instead leaving it released at a logic high level. The slave might issue a NACK for a variety of reasons:
\begin{itemize}
	\item No device exists on the bus with the desired slave address
	\item The slave is busy performing some other task and is presently incapable of communicating with the master
\end{itemize}
If the slave issues a NACK, the master must terminate the transfer by issuing either a stop or restart condition. If the slave issues an ACK, then the transfer continues to the subsequent data frame(s).

\begin{figure}[h]
	\centering
	\begin{tikztimingtable}[timing/font=\rmfamily, timing/d/text/.append style={font=\rmfamily}, xscale=1.5, yscale=2, timing/slope=0.2, timing/dslope=0.2, timing/d/background/.style={fill=white},
		timing/metachar={{K}[2]{#1H !{++(-.5\xunit + 0.5*\slope\xunit, -.5\yunit)} N[rectangle,scale=.6]{#2} !{++(.5\xunit - 0.5*\slope\xunit, +.5\yunit)}}},
		timing/metachar={{J}[2]{#1L !{++(-.5\xunit + 0.5*\slope\xunit, +.5\yunit)} N[rectangle,scale=.6]{#2} !{++(.5\xunit - 0.5*\slope\xunit, -.5\yunit)}}},]
		SDA & H K{IDLE} H ;[blue] L J{START} L 3D{A6} 3D{A5} 3D{A4} 3D{A3} 3D{A2} 3D{A1} 3D{A0} 3D{R/W} ; ;[red] 3D{N/ACK} ; ;[blue] L ; ;[blue, dotted] L ; \\ %H K{STOP} H \\
		SCL & 4H ;[blue] 3L 9{H 2L} ; ;[blue, dotted] L \\ %3H \\
		\begin{extracode}
			\timingaxis{2}
		\end{extracode}
	\end{tikztimingtable}
	\textit{Black}: Idle,~~~{\color{blue}\textit{Blue}: Master-issued},~~~{\color{red}\textit{Red}: Slave-issued}
	\caption{I2C Address Frame Timing Diagram}
	\label{fig:i2c-address-frame-timing-diagram}
\end{figure}

\subsubsection{Data Frame}

In the data frame, the transmitting device sends an 8-bit byte of data to the receiving device. The data frame can be in one of two modes:
\begin{itemize}
	\item Data is sent from the master to the slave. From the master's point of view, it is in \textit{master transmitter} mode. From the slave's point of view, it is in \textit{slave receiver} mode.
	\item Data is sent from the slave to the master. From the master's point of view, it is in \textit{master receiver} mode. From the slave's point of view, it is in \textit{slave transmitter} mode.
\end{itemize}

In either case, the data frame consists of 8 data bits issued by the transmitter followed by an ACK/NACK bit issued by the receiver. The master is always responsible for issuing clock pulses on SCL, regardless of the mode. This is true even in master receiver mode, where the slave is controlling SDA, but must do so in synchronization with the master's SCL pulses. Figure~\ref{fig:i2c-data-frame-timing-diagram-master-transmitter} shows the data frame in master transmitter/slave receiver mode, whereas Figure~\ref{fig:i2c-data-frame-timing-diagram-master-receiver} shows the data frame in master receiver/slave transmitter mode. Note that the only difference between the two is which device is controlling SDA for each bit.

\begin{figure}[h]
	\centering
	\begin{subfigure}[t]{\textwidth}
		\centering
		\begin{tikztimingtable}[timing/font=\rmfamily, timing/d/text/.append style={font=\rmfamily}, xscale=1.5, yscale=2, timing/slope=0.2, timing/dslope=0.2, timing/d/background/.style={fill=white},
			timing/metachar={{K}[2]{#1H !{++(-.5\xunit + 0.5*\slope\xunit, -.5\yunit)} N[rectangle,scale=.6]{#2} !{++(.5\xunit - 0.5*\slope\xunit, +.5\yunit)}}},
			timing/metachar={{J}[2]{#1L !{++(-.5\xunit + 0.5*\slope\xunit, +.5\yunit)} N[rectangle,scale=.6]{#2} !{++(.5\xunit - 0.5*\slope\xunit, -.5\yunit)}}},]
			SDA & ;[dotted] D{} ; ;[blue] 3D{D7} 3D{D6} 3D{D5} 3D{D4} 3D{D3} 3D{D2} 3D{D1} 3D{D0} ; ;[red] 3D{N/ACK} ; ;[blue] L ; ;[blue, dotted] L ; \\
			SCL & ;[blue, dotted] L ;[blue] L 9{H 2L} ; ;[blue, dotted] L \\
			\begin{extracode}
				\timingaxis{2}
			\end{extracode}
		\end{tikztimingtable}
		\caption{Master Transmitter/Slave Receiver}
		\label{fig:i2c-data-frame-timing-diagram-master-transmitter}
	\end{subfigure}
	\par\bigskip\bigskip
	\begin{subfigure}[b]{\textwidth}
		\centering
		\begin{tikztimingtable}[timing/font=\rmfamily, timing/d/text/.append style={font=\rmfamily}, xscale=1.5, yscale=2, timing/slope=0.2, timing/dslope=0.2, timing/d/background/.style={fill=white},
			timing/metachar={{K}[2]{#1H !{++(-.5\xunit + 0.5*\slope\xunit, -.5\yunit)} N[rectangle,scale=.6]{#2} !{++(.5\xunit - 0.5*\slope\xunit, +.5\yunit)}}},
			timing/metachar={{J}[2]{#1L !{++(-.5\xunit + 0.5*\slope\xunit, +.5\yunit)} N[rectangle,scale=.6]{#2} !{++(.5\xunit - 0.5*\slope\xunit, -.5\yunit)}}},]
			SDA & ;[dotted] D{} ; ;[red] 3D{D7} 3D{D6} 3D{D5} 3D{D4} 3D{D3} 3D{D2} 3D{D1} 3D{D0} ; ;[blue] 3D{N/ACK} ; L ;[dotted] L ; \\
			SCL & ;[blue, dotted] L ;[blue] L 9{H 2L} ; ;[blue, dotted] L \\
			\begin{extracode}
				\timingaxis{2}
			\end{extracode}
		\end{tikztimingtable}
		\caption{Master Receiver/Slave Transmitter}
		\label{fig:i2c-data-frame-timing-diagram-master-receiver}
	\end{subfigure}

	{\color{blue}\textit{Blue}: Master-issued},~~~{\color{red}\textit{Red}: Slave-issued},~~~\textit{Black}: Undetermined
	\caption{I2C Data Frame Timing Diagram}
	\label{fig:i2c-data-frame-timing-diagram}
\end{figure}

\subsubsection{Master Transmitter/Slave Receiver Mode}

If the master is the transmitter, it issues the 8 data bits on SDA beginning with the MSB and the slave receives them. The slave latches each consecutive bit on the rising edges of SCL, which are issued by the master. After the 8 data bits, the slave then seizes control of SDA to send the ACK/NACK bit. It will send an ACK if it received the data byte correctly and is ready for another. However, there are a variety of reasons it may send a NACK instead:
\begin{itemize}
	\item The slave is busy performing some other task and is presently incapable of communicating with the master
	\item The slave received data that it considers invalid or incomprehensible
	\item The slave cannot receive any more data bytes (this is the normal way that the slave tells the master to stop sending it more data, and does not necessarily indicate an error)
\end{itemize}

After the data frame, the master may do one of three actions: proceed to send the slave another data byte, issue a stop condition and release control of the bus, or issue a restart condition to retain control of the bus but begin a new transmission (perhaps with another slave, or with the same slave but in master receiver/slave transmitter mode).

\subsubsection{Master Receiver/Slave Transmitter Mode}

If the master is the receiver, the slave issues the 8 data bits on SDA beginning with the MSB and the master receives them. The master latches each consecutive bit on the rising edges of SCL, which are issued (as always) by the master. After the 8 data bits, the master then seizes control of SDA to send the ACK/NACK bit.

The master sends an ACK if it wants the slave to continue sending it more data frames. It would then provide more SCK pulses to make the slave send another data frame. Conversely, the master sends a NACK if it wants to conclude the transfer. It must then issue a stop to release control of the bus, or restart condition to retain control of the bus but begin a new transmission (perhaps with another slave, or with the same slave but in master transmitter/slave receiver mode).



\subsection{Clock Stretching}

In some situations, the slave may be unable to reply to the master because the master's SCL clock rate and time between data frames is too fast for the slave to process. Ordinarily, the slave would simply issue a NACK to tell the master that is was not able to proceses the data in the transfer. This is undesirable because it does not convey to the master why the slave issued a NACK, and the master may not realize that the solution is to slow down. In such a situation, the master might keep trying and failing to communicate the slave over and over again.

Clock stretching eliminates this problem by allowing the slave to pause the current transfer until it is ready to continue. It does this by momentarily seizing control of SCL by driving it low. This prevents the master from issuing any further clock pulses on SCL until the slave permits it because the master cannot make a rising edge on SCL until the slave releases it from being driven low.

On this chip, slave clock strecthing is enabled when \bitfield{I2CSCS} is 1. In this mode, the slave will automatically drive the SCL wire low during the ACK/NACK bit until the slave continue bit \bitfield{I2CSC} in the I2C flow control register \register{I2CxFCL} is set to 1, whereupon the slave will immediately release SCL and allow the master to continue issuing clock pulses. The user is tasked with doing its processing while the clock strectching is engaged, and then setting the ACK/NACK bit and then continuing the transfer.

Not all I2C masters support slaves with clock stretching, so the user must investigate if the desired external device supports it or not.



\subsection{Arbitration}

In situations where more than one device on the bus could possibly become the master, it is possible that two or more devices will attempt to issue a start condition and seize control of the bus at the exact same moment. The I2C bus protocol is not capable of having two masters control the bus at the same time, but the competing masters would not realize their error because multiple masters issuing a start condition simultaneously is indistinguishable from only a single master doing so. Therefore, the I2C protocol provides a process called arbitration to decide which of the competing masters will actually gain control of the bus, while the others release control and cease to claim to be the master.

To perform arbitration, each claimant master attempts to send each bit in the address frame and monitors the state of SDA when SCL has a rising edge to check if the value of SDA matches the value it is attempting to send on it. If the measured value does not match the attempted value, the device knows that some other device is asserting control over SDA. Therefore, the master that noticed the value mismatch knows that it has lost arbitration, and immediately releases control of the bus and ceases to participate in the transfer. Note that whichever master is driving SDA low is the one that will win arbitration, and any that are trying to let SDA be high will lose. The master that wins will not ever realize any other device was trying to seize control of the bus. If all masters are sending the same bit, they proceed on to the next bit until a difference is detected. If all claimant masters attempt to contact the same slave (i.e. the address frame is identical), arbitration will proceed to the data frame(s). If all claimant masters send exactly the same message, then no error occurs because the message was sent successfully.

No information is lost during the arbitration process. If a device loses arbitration, it must wait for the bus to be released before attempting to seize the bus again. It is possible that the winning master is attempting to address the losing device as a slave, in which case that device must be ready to become a slave as soon as it detects an arbitration loss. There is no in-built priority to determining which master will win arbitration.

On this chip, if the I2C peripheral is set up as a master and this device loses arbitration, the \bitfield{I2CMARB} flag will be set. If the \bitfield{I2CMARBIE} bit is set, an interrupt will be generated.



\subsection{Reserved Addresses}

There are several addresses that are not allowed for a slave to have in the I2C protocol. They each have specific purposes:

\begin{itemize}
	\item Address = \texttt{0000 000}, R/W bit = \texttt{0}: General call address
	\item Address = \texttt{0000 000}, R/W bit = \texttt{1}: START byte
	\item Address = \texttt{0000 001}, R/W bit = \texttt{X}: CBUS address
	\item Address = \texttt{0000 010}, R/W bit = \texttt{X}: Reserved for different bus format
	\item Address = \texttt{0000 011}, R/W bit = \texttt{X}: Reserved for future purposes
	\item Address = \texttt{0000 1XX}, R/W bit = \texttt{X}: Hs-mode controller code
	\item Address = \texttt{1111 1XX}, R/W bit = \texttt{1}: Device ID
	\item Address = \texttt{1111 0XX}, R/W bit = \texttt{X}: 10-bit target addressing
\end{itemize}



\subsection{I2C General Call}

If a master wishes to transmit data to every potential slave on the I2C bus all at once, it can issue a general call transfer by setting the desired slave address to \texttt{0000 000} and the R/W bit to \texttt{0} for write mode in the address frame. At this point, any potential slave that wishes to listen for a general call will enter slave receiver mode and send an ACK. Note that if no slaves send an ACK, the master will know that no device is listening. However, if at least one slave issues an ACK, then the master will know that at least one slave is listening, but will not know how many. The same is the case for the following data frames: if one slave sends an ACK, the master will not know how many (if any) are sending a NACK.

On this chip, the ability to listen for a general call as a slave can be enabled or disabled with the \bitfield{I2CGCE} bit.



\subsection{I2C Baud Clock}

When in master mode, the frequency of the \pin{SCLx} clock pin is determined by the frequency of SMCLK and the I2C clock divider bits \bitfield{I2CMDIV}, given by the expression:
\begin{equation*}
f_\textrm{SCL} = \frac{f_\textrm{SMCLK}}{4 \times 2^\bitfield{I2CMDIV}}
\end{equation*}



\subsection{Transmit Buffer Register Queue}

Writing to the I2C transmit buffer registers \register{I2CxMTX} or \register{I2CxSTX} queues up the word written to it as the next frame to transmit. If a transfer is not currently ongoing, the I2C peripheral begins transmitting the word written to the appropriate transmit register as soon as possible, and the corresponding transmit register empty bit is then set. If a transfer is currently ongoing, the I2C peripheral waits for the current frame to finish transferring. As soon as it finishes, the I2C peripheral begins transmitting the new frame written to the appropriate transmit register and sets the corresponding transmit register empty bit.

The TX queue allows the I2C peripheral to constantly transfer data without any dead time in between frames, ensuring maximum transfer speed. However, the user must take care not to overflow the queue, which happens if the appropriate transmit register is written two or more times during a single frame. Queue overflows can be prevented by waiting for the corresponding transmit register empty bit to be 1 before writing to the transmit register.



\subsection{Flags and Interrupts}

The \peripheral{I2Cx} peripheral has 13 interrupt flags and 3 non-interrupt indicators. If the corresponding interrupts are enabled in the \register{SPIxCR} register and the \peripheral{I2Cx} ISR is not masked, then the interrupt flags will generate an interrupt whenever any one of them is equal to 1. The interrupt flags are cleared by writing anything to the status register \register{I2CxSR}, which must be done before the \peripheral{I2Cx} ISR returns, else the ISR will mistakenly re-execute immediately upon the prior return. The interrupt flags are:

\begin{itemize}
	\item \bitfield{I2CSPR}: Stop condition received interrupt flag. This flag is set whenever a stop condition condition is detected on the bus, regardless of which device sent it.
	\item \bitfield{I2CSTR}: Start (or restart) condition received interrupt flag. This flag is set whenever a start condition or repeated start condition is detected on the bus, regardless of which device sent it.
	\item \bitfield{I2CMXC}: Master transfer complete interrupt flag. In master transmitter mode or during the address frame of master receiver mode, this flag is set after this master has sent the data byte and the slave has sent an ACK/NACK. In master receiver mode, this flag is set after the slave has sent the data byte and before this master sends an ACK/NACK.
	\item \bitfield{I2CMNR}: Master mode NACK received interrupt flag. This flag is set in master transmitter mode or during the address frame of master receiver mode after ACK/NACK bit is sent if the slave sends a NACK. It is not set if the slave sends an ACK.
	\item \bitfield{I2CMTXE}: Master transmit register empty interrupt flag. This bit is set when this master latches the data stored in the master transmit register to indicate that the master transmit register is ready to accept another byte and queue it for transmission.
	\item \bitfield{I2CMARB}: Master mode arbitration loss interrupt flag. This bit is set when this master detects that the value it tried to write to SDA is being overridden by another master. After it detects the arbitration loss, this master immediately releases control of the bus.
	\item \bitfield{I2CMSPS}: Master mode stop condition sent interrupt flag. This flag is set after this master sends a stop condition. Note that the \bitfield{I2CSPR} flag is set whenever \textit{any} stop condition is detected, including when this master sends one.
	\item \bitfield{I2CMSTS}: Master mode start condition sent interrupt flag. This flag is set after this master sends a start condition or repeated start condition. Note that the \bitfield{I2CSTR} flag is set whenever \textit{any} start/restart condition is detected, including when this master sends one.
	\item \bitfield{I2CSXC}: Slave mode transfer complete interrupt flag. This bit is set after this slave receives a byte of data from a master, but before this slave sends an ACK/NACK.
	\item \bitfield{I2CSNR}: Slave mode NACK received from master interrupt flag. This bit is set in slave transmitter mode if the master responds with a NACK after this slave sends it a byte of data. This bit is not set if the master responds with an ACK.
	\item \bitfield{I2CSOVF}: Slave receive register overflow interrupt flag. Indicates that this slave has failed to read one or more bytes from the I2CxSRX register before they were overwritten by another transmission.
	\item \bitfield{I2CSTXE}: Slave transmit register empty interrupt flag. This bit is set when this slave latches the data stored in the slave transmit register to indicate that the slave transmit register is ready to accept another byte and queue it for transmission.
	\item \bitfield{I2CSA}: Slave mode addressed interrupt flag. Indicates that this device has been addressed by another master during the address frame, and that it has been designated as the slave for the duration of the transfer.
	\item \bitfield{I2CSTM}: Slave transmitter mode non-interrupt indicator. Indicates for which mode this slave has been addressed: slave receiver (when 0) or slave transmitter (when 1). Only valid if I2CSA is 1. This bit cannot be cleared by writing to the status register.
	\item \bitfield{I2CMCB}: Master controls bus non-interrupt indicator. Shows if this master has control of the I2C bus or not. This bit cannot be cleared by writing to the status register.
	\item \bitfield{I2CBS}: I2C bus state indicator. Shows if the I2C bus is idle or active. This bit cannot be cleared by writing to the status register.
\end{itemize}



\subsection{Master Transmitter Mode Transfer Procedure}

\begin{figure}[h]
	\centering
	\begin{tikztimingtable}[timing/font=\rmfamily, timing/d/text/.append style={font=\rmfamily}, xscale=1.5, yscale=2, timing/slope=0.2, timing/dslope=0.2, timing/d/background/.style={fill=white},
		timing/metachar={{K}[2]{#1H !{++(-.5\xunit + 0.5*\slope\xunit, -.5\yunit)} N[rectangle,scale=.6]{#2} !{++(.5\xunit - 0.5*\slope\xunit, +.5\yunit)}}},
		timing/metachar={{J}[2]{#1L !{++(-.5\xunit + 0.5*\slope\xunit, +.5\yunit)} N[rectangle,scale=.6]{#2} !{++(.5\xunit - 0.5*\slope\xunit, -.5\yunit)}}},]
		SDA & H K{IDLE} H ;[blue] L J{START} L 3D{A6} ; ;[blue, dotted] D{} ; ;[blue] 3D{A0} 3D{R/W} ; ;[red] 3D{N/ACK} ; ;[blue] L ;    ;[blue, dotted] L ; ;[blue] L 3D{D7} ; ;[blue, dotted] D{} ; ;[blue] 3D{D0} ; ;[red] 3D{N/ACK} ; ;[blue, dotted] L ; ;[blue] 2L ; H K{STOP} H \\
		SCL & 4H ;[blue] 2L L H L ; ;[blue, dotted] L ; ;[blue] 3{L H L} 3L    L H L ; ;[blue, dotted] L ; ;[blue] 2{L H L} ; ;[blue, dotted] L ; ;[blue] L ; 4H \\
		\begin{extracode}
			\timingaxis{2}
		\end{extracode}
	\end{tikztimingtable}
	\\
	\textit{Black}: Idle,~~~{\color{blue}\textit{Blue}: Master-issued},~~~{\color{red}\textit{Red}: Slave-issued}
	\caption{I2C Master Transmitter/Slave Receiver Timing Diagram}
	\label{fig:i2c-master-transmitter}
\end{figure}

The procedure for using the \peripheral{I2Cx} peripheral in master transmitter mode is as follows (see Figure~\ref{fig:i2c-master-transmitter} for more details):

\begin{enumerate}
	\item Set up the \peripheral{I2Cx} peripheral by enabling it in master mode (set \bitfield{I2CMEN} to 1) and setting the master clock divider \bitfield{I2CMDIV} to the desired SCL clock frequency. Clear the status register.
	\item Wait until the I2C bus is idle ($t = 0$). The bus is idle when \bitfield{I2CBS} is 0. If this is the first transfer after the \peripheral{I2Cx} peripheral was enabled, poll \pin{SDAx} and \pin{SCLx} for several milliseconds to make sure the bus is really idle, as the \bitfield{I2CBS} bit does not monitor the I2C bus while the peripheral is disabled.
	\item \label{proc:i2c-master-transmitter-send-start} Send a start condition by writing a 1 to \bitfield{I2CMST} ($t = 2$). The start condition will begin to be sent on the next clock cycle ($t = 3$), which is also when the start condition received flag (\bitfield{I2CSTR}) will become set to 1. The master controls this bus bit \bitfield{I2CMCB} will now be set to 1.
	\item Wait for the master start condition sent flag (\bitfield{I2CMSTS}) to become 1, indicating that the start condition has finished being sent ($t = 4$). Clear the status register.
	\item Write the address frame byte to the master transmit register \register{I2CxMTX} ($t = 5$). In master transmitter mode, the address frame byte has the 7-bit desired slave address in the most significant 7 bits, and a 0 in the least significant bit (the read/write bit, which is a 0 to indicate write, or master transmitter, mode). The \peripheral{I2Cx} peripheral will begin transmitting the address frame on the next clock cycle ($t = 6$).
	\item Wait for the master transfer complete flag \bitfield{I2CMXC} or the master arbitration lost flag \bitfield{I2CMARB} flag to be set. If \bitfield{I2CMXC} is set, then the transfer (including the ACK/NACK bit) has completed ($t = 17$). If \bitfield{I2CMARB} is set, then this master has lost arbitration and has released control of the I2C bus (\bitfield{I2CMCB} will now be set to 0).
	\item Check the master NACK received bit \bitfield{I2CMNR}. If it is a 0, then the desired slave has sent an ACK to indicate that it is ready for the transfer to continue. If it is a 1, then the desired slave is not ready, and this master must either send a stop or a restart condition to terminate the current transfer. Clear the status register.
	\item \label{proc:i2c-master-transmitter-send-data-frame} Send the slave a byte of data in the data frame by writing that byte to \register{I2CxMTX} ($t = 21$). The \peripheral{I2Cx} peripheral will begin transmitting the address frame on the next clock cycle ($t = 22$).
	\item Wait for the master transfer complete flag \bitfield{I2CMXC} or the master arbitration lost flag \bitfield{I2CMARB} flag to be set. If \bitfield{I2CMXC} is set, then the transfer (including the ACK/NACK bit) has completed ($t = 17$). If \bitfield{I2CMARB} is set, then this master has lost arbitration and has released control of the I2C bus (\bitfield{I2CMCB} will now be set to 0).
	\item Check the master NACK received bit \bitfield{I2CMNR}. If it is a 0, then the desired slave has sent an ACK to indicate that it has successfully received the data. If it is a 1, then the desired slave has encountered a problem or cannot receive any more data, and this master must either send a stop or a restart condition to terminate the current transfer. Clear the status register.
	\item Continue sending more data bytes (loop back to step \ref{proc:i2c-master-transmitter-send-data-frame}), or terminate the current transfer by sending a restart condition (loop back to step \ref{proc:i2c-master-transmitter-send-start}) or a stop condition (by writing a 1 to \bitfield{I2CMSP} at $t = 33$. The stop condition will begin to be sent at $t = 34$, which is also when the stop condition received flag \bitfield{I2CSPR} will become set to 1).
	\item Wait for the master stop condition sent flag (\bitfield{I2CMSPS}) to become 1, indicating that the stop condition has finished being sent ($t = 35$). Clear the status register.
	\item The bus is now idle, and any device can seize control of it and become the new master.
\end{enumerate}



\subsection{Master Receiver Mode Transfer Procedure}

\begin{figure}[h]
	\centering
	\begin{tikztimingtable}[timing/font=\rmfamily, timing/d/text/.append style={font=\rmfamily}, xscale=1.5, yscale=2, timing/slope=0.2, timing/dslope=0.2, timing/d/background/.style={fill=white},
		timing/metachar={{K}[2]{#1H !{++(-.5\xunit + 0.5*\slope\xunit, -.5\yunit)} N[rectangle,scale=.6]{#2} !{++(.5\xunit - 0.5*\slope\xunit, +.5\yunit)}}},
		timing/metachar={{J}[2]{#1L !{++(-.5\xunit + 0.5*\slope\xunit, +.5\yunit)} N[rectangle,scale=.6]{#2} !{++(.5\xunit - 0.5*\slope\xunit, -.5\yunit)}}},]
		SDA & H K{IDLE} H ;[blue] L J{START} L 3D{A6} ; ;[blue, dotted] D{} ; ;[blue] 3D{A0} 3D{R/W} ; ;[red] 3D{N/ACK} ; ;[blue] L ;    ;[blue, dotted] L ; ;[blue] L ; ;[red] 3D{D7} ; ;[red, dotted] D{} ; ;[red] 3D{D0} ; ;[blue] 3D{N/ACK} ; ;[blue, dotted] L ; ;[blue] 2L ; H K{STOP} H \\
		SCL & 4H ;[blue] 2L L H L ; ;[blue, dotted] L ; ;[blue] 3{L H L} 3L    L H L ; ;[blue, dotted] L ; ;[blue] 2{L H L} ; ;[blue, dotted] L ; ;[blue] L ; 4H \\
		\begin{extracode}
			\timingaxis{2}
		\end{extracode}
	\end{tikztimingtable}
	\\
	\textit{Black}: Idle,~~~{\color{blue}\textit{Blue}: Master-issued},~~~{\color{red}\textit{Red}: Slave-issued}
	\caption{I2C Master Receiver/Slave Transmitter Timing Diagram}
	\label{fig:i2c-master-receiver}
\end{figure}

The procedure for using the \peripheral{I2Cx} peripheral in master receiver mode is as follows (see Figure~\ref{fig:i2c-master-receiver} for more details):

\begin{enumerate}
	\item Set up the \peripheral{I2Cx} peripheral by enabling it in master mode (set \bitfield{I2CMEN} to 1) and setting the master clock divider \bitfield{I2CMDIV} to the desired SCL clock frequency. Clear the status register.
	\item Wait until the I2C bus is idle ($t = 0$). The bus is idle when \bitfield{I2CBS} is 0. If this is the first transfer after the \peripheral{I2Cx} peripheral was enabled, poll \pin{SDAx} and \pin{SCLx} for several milliseconds to make sure the bus is really idle, as the \bitfield{I2CBS} bit does not monitor the I2C bus while the peripheral is disabled.
	\item \label{proc:i2c-master-receiver-send-start} Send a start condition by writing a 1 to \bitfield{I2CMST} ($t = 2$). The start condition will begin to be sent on the next clock cycle ($t = 3$), which is also when the start condition received flag (\bitfield{I2CSTR}) will become set to 1. The master controls this bus bit \bitfield{I2CMCB} will now be set to 1.
	\item Wait for the master start condition sent flag (\bitfield{I2CMSTS}) to become 1, indicating that the start condition has finished being sent ($t = 4$). Clear the status register.
	\item Write the address frame byte to the master transmit register \register{I2CxMTX} ($t = 5$). In master receiver mode, the address frame byte has the 7-bit desired slave address in the most significant 7 bits, and a 1 in the least significant bit (the read/write bit, which is a 1 to indicate read, or master receiver, mode). The \peripheral{I2Cx} peripheral will begin transmitting the address frame on the next clock cycle ($t = 6$).
	\item Wait for the master transfer complete flag \bitfield{I2CMXC} or the master arbitration lost flag \bitfield{I2CMARB} flag to be set. If \bitfield{I2CMXC} is set, then the transfer (including the ACK/NACK bit) has completed ($t = 17$). If \bitfield{I2CMARB} is set, then this master has lost arbitration and has released control of the I2C bus (\bitfield{I2CMCB} will now be set to 0).
	\item Check the master NACK received bit \bitfield{I2CMNR}. If it is a 0, then the desired slave has sent an ACK to indicate that it is ready for the transfer to continue. If it is a 1, then the desired slave is not ready, and this master must either send a stop or a restart condition to terminate the current transfer. Clear the status register.
	\item Receive a data byte from the slave by writing a 1 to the master receive byte bit \bitfield{I2CMRB} ($t = 21$). The master will begin issuing clock pulses on SCL ($t = 22$), and the slave will begin transmitting data on SDA simultaneously.
	\item Wait for the master transfer complete flag \bitfield{I2CMXC} to be set ($t = 28$), indicating that the byte of data from the slave has been received and is ready to be read in the master receive register \register{I2CxMRX}. Read the data in \register{I2CxMRX}. Note that at this point, the master has not sent the ACK/NACK bit yet. Clear the status register.
	\item If the master wishes to receive another byte of data, write a 1 to \bitfield{I2CMRB}, which will immediately send an ACK and then begin another data frame and receive another byte. Continue doing this until no more data bytes are desired.
	\item Terminate the current transfer by either sending a restart condition (loop back to step \ref{proc:i2c-master-receiver-send-start}) or a stop condition (by writing a 1 to \bitfield{I2CMSP} at $t = 28$). In either case, the master will first issue a NACK, and then proceed to issue the stop/restart condition. The stop condition will begin to be sent at $t = 34$, which is also when the stop condition received flag \bitfield{I2CSPR} will become set to 1.
	\item Wait for the master stop condition sent flag (\bitfield{I2CMSPS}) to become 1, indicating that the stop condition has finished being sent ($t = 35$). Clear the status register.
	\item The bus is now idle, and any device can seize control of it and become the new master.
\end{enumerate}



\subsection{Slave Transmitter Mode Transfer Procedure}

The procedure for using the \peripheral{I2Cx} peripheral in slave transmitter mode is as follows (see Figure~\ref{fig:i2c-master-receiver} for more details):

\begin{enumerate}
	\item Set up the \peripheral{I2Cx} peripheral by setting this slave's 7-bit address (with \register{I2CxAR}) and address mask (with \register{I2CxAMR}), setting up an ACK for the address frame (write a 0 to \bitfield{I2CSN}), setting the clock streching mode (with \bitfield{I2CSCS}) and general call sensitivity (with \bitfield{I2CGCE}). Clear the status register. Then, enable slave mode (write a 1 to \bitfield{I2SEN}).
	\item Wait for the ``this slave addressed'' flag \bitfield{I2CSA} to be set to 1 ($t = 15$), indicating that an external master has addressed this device to become the slave.
	\item If clock stretching is enabled, determine if this slave should ACK or NACK by configuring \bitfield{I2CSN}. Then, queue the first byte to be transmitted to the master by writing that byte to the slave transmit register \register{I2CxSTX}. Then, release SCL and allow the transfer to continue by writing a 1 to the slave continue bit \bitfield{I2CSC} ($t = 30$). If clock streching is not enabled, then this slave will automatically ACK or NACK depending on the predetermined value of \bitfield{I2CSN}, and it is not necessary to write a 1 to \bitfield{I2CSC}.
	\item Check if this transfer is in slave receiver or slave transmitter mode by reading \bitfield{I2CSTM} (1 indicates slave transmitter). Clear the status register.
	\item \label{proc:i2c-slave-transmitter-write} Wait for the slave transmit register to be empty when \bitfield{I2CSTXE} is a 1 ($t = 18$ or $t = 31$), then queue the next byte to transmit. If there are no bytes left to transmit, send a NACK at the end of the current transfer by writing a 1 to \bitfield{I2CSN}.
	\item \label{proc:i2c-slave-transmitter-wait} Wait for the slave transfer complete flag \bitfield{I2CSXC} to be 1 ($t = 28$), or for the transmission to be terminated by waiting for either the stop received flag \bitfield{I2CSPR} ($t = 34$) or the start/restart received flag \bitfield{I2CSTR} to be set to 1.
	\item If clock stretching is enabled, determine if this slave should ACK or NACK by configuring \bitfield{I2CSN}. Then, release SCL and allow the transfer to continue by writing a 1 to the slave continue bit \bitfield{I2CSC} ($t = 30$). If clock streching is not enabled, then this slave will automatically ACK or NACK depending on the predetermined value of \bitfield{I2CSN}, and it is not necessary to write a 1 to \bitfield{I2CSC}.
	\item Read the status register, then clear it. If this was the last byte to transmit, go to step \ref{proc:i2c-slave-transmitter-end}. Otherwise, if \bitfield{I2CSXC} was set, go back to step \ref{proc:i2c-slave-transmitter-write}. If either \bitfield{I2CSPR} or \bitfield{I2CSTR} were set, go to step \ref{proc:i2c-slave-transmitter-end}.
	\item \label{proc:i2c-slave-transmitter-end} Clear the status register. The transfer has ended, and this slave is ready for the next one.
\end{enumerate}



\subsection{Slave Receiver Mode Transfer Procedure}

The procedure for using the \peripheral{I2Cx} peripheral in slave receiver mode is as follows (see Figure~\ref{fig:i2c-master-transmitter} for more details):

\begin{enumerate}
	\item Set up the \peripheral{I2Cx} peripheral by setting this slave's 7-bit address (with \register{I2CxAR}) and address mask (with \register{I2CxAMR}), setting up an ACK for the address frame (write a 0 to \bitfield{I2CSN}), setting the clock streching mode (with \bitfield{I2CSCS}) and general call sensitivity (with \bitfield{I2CGCE}). Clear the status register. Then, enable slave mode (write a 1 to \bitfield{I2SEN}).
	\item Wait for the ``this slave addressed'' flag \bitfield{I2CSA} to be set to 1 ($t = 15$), indicating that an external master has addressed this device to become the slave.
	\item If clock stretching is enabled, determine if this slave should ACK or NACK by configuring \bitfield{I2CSN}. Then, release SCL and allow the transfer to continue by writing a 1 to the slave continue bit \bitfield{I2CSC} ($t = 30$). If clock streching is not enabled, then this slave will automatically ACK or NACK depending on the predetermined value of \bitfield{I2CSN}, and it is not necessary to write a 1 to \bitfield{I2CSC}.
	\item Check if this transfer is in slave receiver or slave transmitter mode by reading \bitfield{I2CSTM} (0 indicates slave receiver). Clear the status register.
	\item \label{proc:i2c-slave-receiver-wait} Wait for either the slave to receive a data byte from the master by waiting for the slave transfer complete flag \bitfield{I2CSXC} to be 1 ($t = 28$), or for the transmission to be terminated by waiting for either the stop received flag \bitfield{I2CSPR} ($t = 34$) or the start/restart received flag \bitfield{I2CSTR} to be set to 1. If \bitfield{I2CSXC} was set, go on to step \ref{proc:i2c-slave-receiver-read}. If either \bitfield{I2CSPR} or \bitfield{I2CSTR} were set, go to step \ref{proc:i2c-slave-transmitter-end}.
	\item \label{proc:i2c-slave-receiver-read} If the slave receive overflow bit \bitfield{I2CSOVF} is a 1, then this slave has failed to read one or more of the bytes sent by the master in the past. Clear the status register. Read the byte sent from the master by reading the slave receive register \register{I2CxSRX}.
	\item If clock stretching is enabled, determine if this slave should ACK or NACK by configuring \bitfield{I2CSN}. Then, release SCL and allow the transfer to continue by writing a 1 to the slave continue bit \bitfield{I2CSC} ($t = 30$). If clock streching is not enabled, then this slave will automatically ACK or NACK depending on the predetermined value of \bitfield{I2CSN}, and it is not necessary to write a 1 to \bitfield{I2CSC}. Note that if this slave cannot transmit any more data bytes, it should NACK after the last data byte.
	\item Loop back to step \ref{proc:i2c-slave-receiver-wait}.
	\item \label{proc:i2c-slave-receiver-end} Clear the status register. The transfer has ended, and this slave is ready for the next one.
\end{enumerate}


